\section{Results}
\label{sec:results}

\subsection{Congruence Varies Across Profiles}
\label{sec:congruence_results}

Not all persona profiles are equally easy for LLMs to adopt.
\Tabref{tab:congruence} summarizes congruence statistics across models.
\mini achieves higher scores overall (mean combined congruence 0.896 vs.\ 0.772), consistent with its greater model capacity, but both models show significant variance across profiles.
A Kolmogorov-Smirnov test rejects the null hypothesis of uniform congruence for both models ($p < 10^{-12}$).

\begin{table}[t]
    \centering
    \caption{Congruence statistics across \nprofiles{} persona profiles. \mini achieves higher scores but both models show significant variance, confirming that some preference combinations are harder to adopt.}
    \begin{tabular}{@{}lcccc@{}}
        \toprule
        \textbf{Model} & \textbf{Behavioral} & \textbf{Coherence} & \textbf{Combined} & \textbf{CV (\%)} \\
        \midrule
        \nano & $0.892 \pm 0.022$ & $0.653 \pm 0.075$ & $0.772 \pm 0.043$ & 5.6 \\
        \mini & $0.987 \pm 0.019$ & $0.805 \pm 0.050$ & $0.896 \pm 0.030$ & 3.3 \\
        \bottomrule
    \end{tabular}
    \label{tab:congruence}
\end{table}

Behavioral and coherence scores correlate moderately ($r = 0.41$ for \nano, $r = 0.33$ for \mini, both $p < 0.001$), indicating that the model's own assessment of preference coherence partially predicts its behavioral consistency, though the two measures capture distinct aspects of persona adoption.

\subsection{Person Space Is Low-Rank}
\label{sec:lowrank}

\Tabref{tab:dimensionality} presents the dimensionality analysis of the co-occurrence matrices.
Both models require only \effrank{} components to explain 80\% of variance---significantly fewer than the 23.4 expected under a random baseline ($p < 0.001$ via permutation test with 1,000 iterations).
The participation ratio confirms this finding: 23.9--24.1 effective dimensions versus 29.6 for random matrices.

\begin{table}[t]
    \centering
    \caption{Effective dimensionality of the preference co-occurrence matrix. Both models show significantly lower dimensionality than random baselines, with identical rank at every threshold. Best (lowest rank) results in \textbf{bold}.}
    \begin{tabular}{@{}lccc@{}}
        \toprule
        \textbf{Metric} & \textbf{\nano} & \textbf{\mini} & \textbf{Random} \\
        \midrule
        Rank for 80\% variance & \textbf{21} & \textbf{21} & $23.4 \pm 0.5$ \\
        Rank for 90\% variance & \textbf{28} & \textbf{28} & --- \\
        Rank for 95\% variance & \textbf{34} & \textbf{34} & --- \\
        Participation ratio & \textbf{23.9} & \textbf{24.1} & $29.6 \pm 0.5$ \\
        \bottomrule
    \end{tabular}
    \label{tab:dimensionality}
\end{table}

The top 10 components capture approximately 55\% of variance, compared to roughly 34\% for random matrices.
This 65\% reduction in effective dimensionality means that the space of coherent personas is substantially more constrained than what random preference combinations would predict.

\subsection{Components Are Partially Interpretable}
\label{sec:components}

We examine the loading patterns of the top principal components.

\para{Component 1} (9.2--9.3\% variance) distinguishes an ``outgoing intellectual'' pole (high loadings on public speaking, international travel, data/statistics) from a ``homebody consumer'' pole (road trips, social media, morning routines).

\para{Component 2} (7.8\% variance) separates ``creative/artistic'' preferences (stand-up comedy, concerts, musical instruments) from ``traditional/team-oriented'' preferences (tradition, working in teams).

\para{Components 3--4} (5.9--6.4\% variance) load on food and lifestyle preferences (hosting, vegetarianism, multitasking) versus analytical preferences (puzzles, coffee, classical music).

\para{Component 5} ($\sim$4.7\% variance) contrasts privacy and minimalism with technology enthusiasm.

These components partially align with known personality dimensions but do not map cleanly onto Big Five traits, consistent with the entanglement findings of \citet{bhandari2026personality}.

\subsection{Person Space Is Consistent Across Models}
\label{sec:crossmodel_results}

This is our strongest finding.
\Tabref{tab:crossmodel} presents the cross-model comparison metrics.
The co-occurrence matrices correlate at $r = 0.96$ ($p < 10^{-15}$), and the Mantel permutation test confirms this is not due to shared marginal distributions ($p < 0.001$).
The Procrustes disparity for the top-5 components is only 0.091, indicating near-perfect alignment after rotation.

\begin{table}[t]
    \centering
    \caption{Cross-model comparison between \nano and \mini. The co-occurrence matrices and principal components are highly aligned, indicating shared person-space structure across model sizes.}
    \begin{tabular}{@{}lcc@{}}
        \toprule
        \textbf{Metric} & \textbf{Value} & \textbf{$p$-value} \\
        \midrule
        Profile congruence correlation & $r = 0.44$ & $5.2 \times 10^{-7}$ \\
        Co-occurrence matrix correlation & $\mathbf{r = 0.96}$ & $< 10^{-15}$ \\
        Mantel permutation test & --- & $< 0.001$ \\
        Procrustes disparity (top-5) & $\mathbf{0.091}$ & --- \\
        Procrustes disparity (top-10) & $\mathbf{0.070}$ & --- \\
        \bottomrule
    \end{tabular}
    \label{tab:crossmodel}
\end{table}

\Tabref{tab:component_matching} shows that the top-5 components match one-to-one between models, with correlations exceeding 0.90 in all cases.
Components 3 and 4 swap in their variance ordering between models, but their content is preserved---suggesting that these dimensions capture genuine structure rather than noise.

\begin{table}[t]
    \centering
    \caption{Cross-correlation of top-5 principal components between \nano and \mini. Each \nano component matches a unique \mini component at $r > 0.90$, with components 3 and 4 swapped in ordering.}
    \begin{tabular}{@{}ccc@{}}
        \toprule
        \textbf{\nano Component} & \textbf{Best-Matching \mini Component} & \textbf{Correlation} \\
        \midrule
        1 & 1 & $\mathbf{0.972}$ \\
        2 & 2 & $\mathbf{0.946}$ \\
        3 & 4 & $\mathbf{0.912}$ \\
        4 & 3 & $\mathbf{0.929}$ \\
        5 & 5 & $\mathbf{0.901}$ \\
        \bottomrule
    \end{tabular}
    \label{tab:component_matching}
\end{table}

\subsection{Polarity Effects}
\label{sec:polarity}

We observe an asymmetric polarity effect.
For \mini, profiles with more ``like'' preferences achieve higher congruence ($r = 0.31$, $p < 0.001$), suggesting a positivity bias---the model more easily adopts positive preferences than negative ones.
\nano shows no such effect ($r = 0.09$, $p = 0.32$).
Domain diversity (number of distinct domains in a profile) does not predict congruence for either model, indicating that the low-rank structure is driven by preference interactions rather than domain coverage.

\subsection{Error Analysis}
\label{sec:errors}

The lowest-congruence profiles tend to contain internally contradictory preferences.
For example, a \nano profile scoring 0.59 combined ``enjoys multitasking'' and ``dreads public speaking'' with ``loves working in teams''---combining team orientation with avoidance of public engagement.
Another low-scoring profile paired ``loves large gatherings'' with ``prefers listening alone,'' contradicting the social-solo dimension.
Conversely, the highest-scoring profiles exhibit stereotypically consistent themes (\eg ``prefers cooperation; dreads public speaking; prefers traveling with companions'' forms a coherent ``quiet, cooperative'' persona).
This pattern confirms that person space has genuine structure: profiles near the low-dimensional manifold are adopted fluently, while those far from it produce inconsistent behavior.
